% !TEX program = xelatex
\documentclass[11pt,a4paper]{article}

\usepackage{fontspec}
\usepackage{polyglossia}
\usepackage{amsmath,amssymb,amsfonts}
\usepackage{xcolor}
\usepackage{graphicx}
\usepackage{geometry}

\geometry{left=2cm,right=2cm,top=2cm,bottom=2cm}

\setdefaultlanguage{persian}
\setotherlanguage{english}
\setmainfont{XB Niloofar}

\newcommand{\en}[1]{\textenglish{#1}}
\newcommand{\separator}{\vspace{1em}\noindent\rule{\linewidth}{0.4pt}\vspace{1em}}
\newcommand{\thickseparator}{\vspace{1em}\noindent\rule{\linewidth}{1pt}\vspace{1em}}

\title{\textbf{پاسخ سوال ۲} \\[0.4cm] \small آموزش رگرسیون لجستیک با \en{Regularization} های \en{L1} و \en{L2}}
\author{محمد باغبان زاده \\ شماره دانشجویی: شناسهٔ دانشجو}
\date{}
\linespread{1.35}

\begin{document}

\maketitle
\thickseparator

\section{مفاهیم پایه: \en{Regularization} در رگرسیون لجستیک (بخش \en{۴-۱۴} کتاب)}

در بخش \textbf{\en{۴-۱۴}} کتاب (\en{Advanced: Regularization})، مفهوم \textbf{منظم‌سازی (\en{Regularization})} به عنوان راهکاری برای جلوگیری از \textbf{بیش‌برازش (\en{overfitting})} معرفی شده است. زمانی که مدل بیش از حد به داده‌های آموزشی منطبق می‌شود، ممکن است نویزهای تصادفی را نیز یاد بگیرد و در نتیجه عملکرد خوبی روی داده‌های تست نداشته باشد. برای حل این مشکل، یک عبارت منظم‌سازی \(R(\theta)\) به تابع هزینه اضافه می‌شود.

بر اساس معادله \en{۴-۴۹} کتاب، تابع هزینه با \en{Regularization} به صورت زیر است (به صورت کمینه‌سازی تابع هزینه منفی لگاریتم درست‌نمایی):

\[
\hat{\theta} = \arg\min_{\theta} \left( -\sum_{i=1}^{m} \log P(y^{(i)}|x^{(i)};\theta) + \alpha R(\theta) \right)
\]

دو نوع رایج منظم‌سازی در کتاب معرفی شده است:

\subsection{منظم‌سازی \en{L2} ( \en{Ridge} )}
طبق معادله \en{۴-۵۰} و \en{۴-۵۱} کتاب:

\[
R(\theta) = \|\theta\|_2^2 = \sum_{j=1}^{n} w_j^2
\]

تابع هزینه منظم‌شده:
\[
\hat{\theta} = \arg\min_{\theta} \left( -\sum_{i=1}^{m} \log P(y^{(i)}|x^{(i)};\theta) + \alpha \sum_{j=1}^{n} w_j^2 \right)
\]

مشتق عبارت منظم‌سازی نسبت به \(w_j\):
\[
\frac{\partial}{\partial w_j} (\alpha w_j^2) = 2\alpha w_j
\]

\newpage
\subsection{منظم‌سازی \en{L1} ( \en{LASSO} )}
طبق معادله \en{۴-۵۲} و \en{۴-۵۳} کتاب:

\[
R(\theta) = \|\theta\|_1 = \sum_{j=1}^{n} |w_j|
\]

تابع هزینه منظم‌شده:
\[
\hat{\theta} = \arg\min_{\theta} \left( -\sum_{i=1}^{m} \log P(y^{(i)}|x^{(i)};\theta) + \alpha \sum_{j=1}^{n} |w_j| \right)
\]

مشتق عبارت منظم‌سازی نسبت به \(w_j\) (به جز در نقطه صفر):
\[
\frac{\partial}{\partial w_j} (\alpha |w_j|) = \alpha \cdot \text{sign}(w_j)
\]

که در آن تابع علامت به صورت زیر تعریف می‌شود:
\[
\text{sign}(w_j) = 
\begin{cases} 
+1 & \text{\en{if }} w_j > 0 \\
-1 & \text{\en{if }} w_j < 0 \\
\text{تعریف‌نشده} & \text{\en{if }} w_j = 0
\end{cases}
\]

\separator
\section{داده‌های مسأله}

\begin{align*}
\text{:بردار ویژگی} &\quad x = [2,\; -1] \\
\text{:برچسب واقعی} &\quad y = 1 \\
\text{:وزن‌های فعلی} &\quad w = [0.5,\; -0.25] \\
\text{:مقدار بایاس} &\quad b = 0 \\
\text{:پارامتر منظم‌سازی} &\quad \lambda = \alpha = 0.1
\end{align*}

\separator
\section{گام ۱: محاسبه خروجی مدل و خطا}

ابتدا مقدار \(z = w \cdot x + b\) را محاسبه می‌کنیم:

\[
z = (0.5 \times 2) + (-0.25 \times -1) + 0 = 1 + 0.25 = 1.25
\]

سپس خروجی مدل (احتمال کلاس مثبت) را با تابع سیگموئید محاسبه می‌کنیم:

\[
\hat{y} = \sigma(z) = \frac{1}{1 + e^{-1.25}}
\]

مقدار \(e^{-1.25} \approx 0.2865\)، بنابراین:

\[
\hat{y} \approx \frac{1}{1 + 0.2865} = \frac{1}{1.2865} \approx 0.7773
\]

خطای مدل (تفاوت بین خروجی و برچسب واقعی):

\[
\hat{y} - y = 0.7773 - 1 = -0.2227
\]

\separator
\section{گام ۲: محاسبه گرادیان بدون \en{Regularization}}

بر اساس بخش \textbf{\en{۴-۶-۱} کتاب} (معادله \en{۴-۲۶})، گرادیان تابع هزینه کراس-آنتروپی نسبت به هر وزن \(w_j\) به صورت زیر است:

\[
\frac{\partial L}{\partial w_j} = (\hat{y} - y) \cdot x_j
\]

بنابراین:

\[
\frac{\partial L}{\partial w_1} = (\hat{y} - y) \cdot x_1 = (-0.2227) \times 2 = -0.4454
\]

\[
\frac{\partial L}{\partial w_2} = (\hat{y} - y) \cdot x_2 = (-0.2227) \times (-1) = 0.2227
\]

\begin{center}
\fbox{\parbox{0.7\linewidth}{\centering
\textbf{گرادیان بدون \en{Regularization}:} \\
\[
\nabla L = [-0.4454,\; 0.2227]
\]
}}
\end{center}

\separator
\section{گام ۳: محاسبه گرادیان با \en{L2 Regularization} ( \en{Ridge} )}

در منظم‌سازی \en{L2}، عبارت منظم‌سازی به صورت \(\lambda \sum w_j^2\) به تابع هزینه اضافه می‌شود. مشتق این عبارت نسبت به \(w_j\) برابر با \(2\lambda w_j\) است. بنابراین گرادیان کل به صورت زیر محاسبه می‌شود:

\[
\frac{\partial L_{L2}}{\partial w_j} = \frac{\partial L}{\partial w_j} + 2\lambda w_j
\]

\begin{align*}
\frac{\partial L_{L2}}{\partial w_1} &= -0.4454 + 2 \times 0.1 \times 0.5 = -0.4454 + 0.1 = -0.3454 \\[4pt]
\frac{\partial L_{L2}}{\partial w_2} &= 0.2227 + 2 \times 0.1 \times (-0.25) = 0.2227 - 0.05 = 0.1727
\end{align*}

\begin{center}
\fbox{\parbox{0.7\linewidth}{\centering
\textbf{گرادیان با \en{L2 Regularization}:} \\
\[
\nabla L_{L2} = [-0.3454,\; 0.1727]
\]
}}
\end{center}

\separator
\section{گام ۴: محاسبه گرادیان با \en{L1 Regularization} ( \en{LASSO} )}

در منظم‌سازی \en{L1}، عبارت منظم‌سازی به صورت \(\lambda \sum |w_j|\) به تابع هزینه اضافه می‌شود. مشتق این عبارت نسبت به \(w_j\) برابر با \(\lambda \cdot \text{sign}(w_j)\) است. بنابراین گرادیان کل به صورت زیر محاسبه می‌شود:

\[
\frac{\partial L_{L1}}{\partial w_j} = \frac{\partial L}{\partial w_j} + \lambda \cdot \text{sign}(w_j)
\]

برای \(w_1 = 0.5 > 0\)، علامت \(w_1\) برابر \(+1\) است.
برای \(w_2 = -0.25 < 0\)، علامت \(w_2\) برابر \(-1\) است.

\begin{align*}
\frac{\partial L_{L1}}{\partial w_1} &= -0.4454 + 0.1 \times (+1) = -0.4454 + 0.1 = -0.3454 \\[4pt]
\frac{\partial L_{L1}}{\partial w_2} &= 0.2227 + 0.1 \times (-1) = 0.2227 - 0.1 = 0.1227
\end{align*}

\begin{center}
\fbox{\parbox{0.7\linewidth}{\centering
\textbf{گرادیان با \en{L1 Regularization}:} \\
\[
\nabla L_{L1} = [-0.3454,\; 0.1227]
\]
}}
\end{center}

\separator
\section{مقایسه سه گرادیان محاسبه‌شده}

\begin{center}
\begin{tabular}{|c|c|c|c|}
\hline
\textbf{وزن} & \textbf{بدون \en{Reg}} & \textbf{\en{L2 Ridge}} & \textbf{\en{L1 LASSO}} \\
\hline
\(w_1 = 0.5\) & \(-0.4454\) & \(-0.3454\) & \(-0.3454\) \\
\hline
\(w_2 = -0.25\) & \(0.2227\) & \(0.1727\) & \(0.1227\) \\
\hline
\end{tabular}
\end{center}

\subsection{تحلیل مقایسه‌ای}

\begin{enumerate}
    \item \textbf{تأثیر روی \(w_1\) (وزن مثبت):}
    \begin{itemize}
        \item هر دو نوع منظم‌سازی، مقدار گرادیان را \textbf{افزایش} دادند (از \(-0.4454\) به \(-0.3454\)).
        \item دلیل: عبارت منظم‌سازی \(2\lambda w_1\) برای \en{L2} و \(\lambda\) برای \en{L1} هر دو مثبت هستند و به گرادیان منفی اضافه می‌شوند.
        \item در این مثال خاص، هر دو منظم‌سازی مقدار یکسانی برای \(w_1\) تولید کردند.
    \end{itemize}
    
    \item \textbf{تأثیر روی \(w_2\) (وزن منفی):}
    \begin{itemize}
        \item در \en{L2}: گرادیان از \(0.2227\) به \(0.1727\) \textbf{کاهش} یافت (چون \(2\lambda w_2\) منفی بود).
        \item در \en{L1}: گرادیان از \(0.2227\) به \(0.1227\) \textbf{کاهش بیشتری} یافت (چون \(\lambda \cdot \text{sign}(w_2) = -0.1\)).
        \item \textbf{نکته مهم:} در منظم‌سازی \en{L1}، تأثیر روی وزن‌های منفی بیشتر است زیرا علامت \(\text{sign}(w_2) = -1\) است.
    \end{itemize}
    
    \item \textbf{تفاوت کلیدی بین \en{L1} و \en{L2}:}
    \begin{itemize}
        \item مطابق با بخش \textbf{\en{۴-۱۴} کتاب}: منظم‌سازی \en{L2} وزن‌ها را به سمت صفر میل می‌دهد اما به ندرت دقیقاً صفر می‌کند (وزن‌های کوچک اما غیرصفر).
        \item منظم‌سازی \en{L1} می‌تواند وزن‌ها را \textbf{دقیقاً صفر} کند و راه‌حل‌های \textbf{پراکنده (\en{sparse})} ایجاد کند.
        \item در این مثال، \(\nabla L_{L1}\) برای \(w_2\) کوچک‌تر از \(\nabla L_{L2}\) است، که نشان می‌دهد \en{L1} فشار بیشتری برای کاهش وزن‌های منفی اعمال می‌کند.
    \end{itemize}
\end{enumerate}

\separator
\section{جمع‌بندی و پاسخ نهایی}

\begin{center}
\fbox{\parbox{0.95\linewidth}{\centering
\textbf{پاسخ نهایی:}
\begin{itemize}
    \item \textbf{گرادیان بدون \en{Regularization}:} \(\nabla L = [-0.4454,\; 0.2227]\)
    \item \textbf{گرادیان با \en{L2} (\en{Ridge}):} \(\nabla L_{L2} = [-0.3454,\; 0.1727]\)
    \item \textbf{گرادیان با \en{L1} (\en{LASSO}):} \(\nabla L_{L1} = [-0.3454,\; 0.1227]\)
\end{itemize}
\textbf{مقایسه:} هر دو نوع منظم‌سازی باعث کاهش قدر مطلق گرادیان شدند. تأثیر \en{L1} روی وزن منفی \(w_2\) قوی‌تر از \en{L2} بود، که این امر با خاصیت \textbf{پراکندگی (\en{sparsity})} منظم‌سازی \en{L1} همخوانی دارد و در بلندمدت می‌تواند برخی وزن‌ها را دقیقاً به صفر برساند.
}}
\end{center}

\thickseparator

\end{document}